\chapter{The Minimal Viable Document}

Inspired by the Minimum Viable Product, the \textbf{Minimal Viable Document (MVD)} is the smallest document that achieves its purpose: changing the reader's understanding sufficiently to enable action.

Properties of an MVD:

\begin{itemize}
    \item \textbf{Clear audience}: Who will read this, and what do they need?
    \item \textbf{Specific confusion addressed}: What single problem does this solve?
    \item \textbf{Iteratively improvable}: How can this evolve with feedback?
    \item \textbf{Modular}: How does this connect to other documentation?
\end{itemize}

LLMs excel at MVD creation because they can:

\begin{itemize}
    \item Generate hypotheses about reader needs
    \item Identify structural patterns across document types
    \item Suggest standard components (glossaries, examples, edge cases)
    \item Transform rough notes into structured prose
\end{itemize}
